% Prose for this counterexample.  Every region below is delimited by an
% environment that tools/build.py extracts; machine facts (ids, uid, dates,
% urls, enums) live in case.json.  See counterexamples/AGENTS.md.

\cxtitle{Log-concavity of all Derivatives of the Jacobi-theta kernel}

\begin{cxcredits}
\posedby{Coffey and Csordas \citeyearpar{Csordas2014}}
\foundby{GPT-5.5 (Pro)}{2026-02-22}
\formalizedby{S.~Sra}
\auditedby{S.~Sra}
\contributedby{S.~Sra}
\end{cxcredits}

\begin{cxcontext}
The kernel $\Phi$ below is the Jacobi theta kernel whose cosine transform is the Riemann $\xi$-function. A positive function $g$ is \emph{strictly log-concave} when $(g')^2-gg''>0$, so the assertion $J_n(t)>0$ says that the $(n-1)$st derivative $\Phi^{(n-1)}$ is strictly log-concave. %, for every order $n$ and every real $t$.
\end{cxcontext}

% ---------------- result: theta-derivative-log-concavity ----------------

\begin{cxsource}{theta-derivative-log-concavity}
G.~Csordas \citeyearpar{Csordas2014}, Open Problem 4.13, restating Coffey--Csordas Conjecture 2.5
\end{cxsource}

\begin{cxstatement}{theta-derivative-log-concavity}
([5, Conjecture 2.5])  Show that the derivatives of the Jacobi theta function, $\Phi(t)$, are (strictly) log-concave on $\R$; that is, for each $n\in\mathbb N$,
\[J_n(t):=(\Phi^{(n)}(t))^2-\Phi^{(n-1)}(t)\Phi^{(n+1)}(t)>0\quad\text{for }t\in\R.\]
Here, by (4.2) of the same paper, the Jacobi theta function (without the usual factor $4$) is $\Phi(t):=\sum_{n\ge1}\pi n^2\left(2\pi n^2e^{4t}-3\right)\exp\left(5t-\pi n^2e^{4t}\right)$.
\end{cxstatement}

\begin{cxsummary}{theta-derivative-log-concavity}
Coffey--Csordas Conjecture 2.5 / Csordas Problem 4.13 \citeyearpar{Csordas2014}
\end{cxsummary}

\begin{cxcertificate}{theta-derivative-log-concavity}
Rigorous interval evaluation of $J_9(1/50)<0$.
\end{cxcertificate}

\begin{cxrefutation}
\begin{theorem}[\statusfalse; computer-assisted proof]\label{thm:theta}
At $n=9$ and $t=1/50$, one has $J_9(t)<0$.
\end{theorem}
\begin{proof}
Put $u=\pi m^2e^{4t}$ and define integer polynomials
\[
P_0(u)=2u-3,\qquad P_{k+1}(u)=4uP_k'(u)+(5-4u)P_k(u).
\]
Termwise differentiation gives
\[
\Phi^{(k)}(t)=\sum_{m\ge1}\pi m^2e^{5t-u}P_k(u).
\]
Outward-rounded interval summation, with an explicit geometrically bounded tail, yields
\begin{align*}
\Phi^{(8)}(1/50)&\approx 2.95940368409056694\cdot10^8,\\
\Phi^{(9)}(1/50)&\approx 5.80442287530267756\cdot10^9,\\
\Phi^{(10)}(1/50)&\approx 2.02040519425950258\cdot10^{11},
\end{align*}
and an interval for $J_9(1/50)$ contained strictly in the negative half-line, centered at
\[
-2.61006208371358922\cdot10^{19}.
\]
The source package contains a high-precision evaluator and a Sage \texttt{RealIntervalField} certificate.  The negative margin is many orders of magnitude larger than the certified truncation error.
\end{proof}

\begin{remark}
  GPT-5.5 managed to prove the Tur\'an inequalities for $\Phi^{(n)}$ for $n=2$ and $n=3$; while we were attempting an inductive proof to generalize those proofs to all $n$ is when the counterexample popped up. However, we have not yet verified whether the Tur\'an inequalities hold for all $n \le 8$.
\end{remark}
\end{cxrefutation}


%%% Local Variables:
%%% mode: LaTeX
%%% TeX-master: "../../tex/main"
%%% End:
